% SLOP8412 optional final-report starter. Compile twice with pdfLaTeX.
% Replace every prompt and placeholder. The assessment brief is authoritative.
\documentclass[11pt,a4paper]{article}
\usepackage[T1]{fontenc}
\usepackage[utf8]{inputenc}
\usepackage{lmodern}
\usepackage[a4paper,margin=23mm,headheight=28pt,headsep=14pt]{geometry}
\usepackage{graphicx,xcolor,amsmath,amssymb,booktabs,tabularx}
\usepackage{fancyhdr,titlesec,enumitem,needspace}
\usepackage{hyperref}
\definecolor{SlopGold}{HTML}{996514}
\definecolor{PromptGrey}{HTML}{505050}
\definecolor{RuleGold}{HTML}{B97D1C}
\hypersetup{colorlinks=true,urlcolor=SlopGold,linkcolor=SlopGold,citecolor=SlopGold,
  pdftitle={SLOP8412 Final Report Template},pdfauthor={Slop University}}

% Replace these fields. A compact title block counts within the report pages.
\newcommand{\reporttitle}{[Your report title]}
\newcommand{\studentname}{[Your name]}
\newcommand{\studentid}{[Student ID]}
\newcommand{\reportdate}{[Submission date]}
% Replace this URL with the deployed course URL if your copy lives elsewhere.
\newcommand{\courseurl}{https://comp4020-agentic-coding-studio.github.io/comp4020-ass2-jnheinrich451-eng}

\pagestyle{fancy}
\fancyhf{}
\fancyhead[L]{\raisebox{-7pt}{\includegraphics[height=25pt]{slop-crest.png}}\hspace{8pt}\textsf{Slop / University}}
\fancyhead[R]{\small\textsf{SLOP8412 / Final report}}
\fancyfoot[L]{\small\textsf{\studentid}}
\fancyfoot[R]{\small\thepage}
\renewcommand{\headrulewidth}{0.4pt}
\renewcommand{\headrule}{\hbox to\headwidth{\color{RuleGold}\leaders\hrule height\headrulewidth\hfill}}
\titleformat{\section}{\Large\sffamily\bfseries\color{SlopGold}}{}{0pt}{}
\titleformat{\subsection}{\normalsize\sffamily\bfseries}{}{0pt}{}
\titlespacing*{\section}{0pt}{14pt}{7pt}
\titlespacing*{\subsection}{0pt}{8pt}{4pt}
\setlength{\parindent}{0pt}
\setlength{\parskip}{4pt}
\setlist[itemize]{leftmargin=*,nosep,topsep=4pt}
\newcommand{\prompt}[1]{{\small\color{PromptGrey}\textit{Writing prompt: #1}\par}}
\newcommand{\blank}{\textemdash}

\begin{document}
{\small\sffamily\color{SlopGold}ADVANCED FR\'ECHET INCEPTION DISTANCE}

{\LARGE\sffamily\bfseries\reporttitle\par}
\vspace{8pt}
\studentname\enspace / \studentid\hfill\reportdate

{\small\textbf{Optional starter, not a completed submission.}
Replace all prompts, empty cells and example objects. Submit at most six pages
plus references, together with a separate archive that scores the bench from a
seed. This title block is part of the report, not an extra cover page.
The \href{\courseurl/assessments/final-report/}{assessment brief} governs the task;
the \href{\courseurl/policies/}{policies page} governs reporting and tool disclosure.
Delete this guidance paragraph before submission.}

\section{Part A. The proposal [25 points]}
\prompt{State what a paper should report in place of FID and which failure it
addresses. Define the proposal precisely enough to implement and evaluate it.
For an existing method, cite its source and identify your contribution.}

\subsection{Definition and intended use}
\prompt{Name the inputs, parameters, output, units and direction of improvement.
Say which comparisons are meaningful. Replace the equation placeholder below.}
\begin{equation}
  \widehat{Q}_{\theta}(\mathcal{X}_R,\mathcal{X}_G)
  = \text{[define your proposed metric here]}.
  \label{eq:proposal}
\end{equation}
\prompt{Explain each symbol in Equation~\ref{eq:proposal}. A reporting protocol
may need a precise procedure rather than a single scalar equation.}

\section{Part B. The evaluation [45 points]}
\prompt{Evaluate the proposal against the course's standard. Address each of the
four questions below, or justify its absence. An empty subsection is not a
justification. Keep measured evidence separate from proposed experiments.}

\subsection{Measurement protocol}
\begin{tabularx}{\linewidth}{@{}lX@{}}
\toprule
Setting & Your specification \\
\midrule
Reference and candidates & [Bench R; candidates A, B and C; construction] \\
Sample counts & [$N_R$, $N_G$, and the sample-size ladder for extrapolation] \\
Implementation & [Code version, estimator, precision and dependencies] \\
Representation & [Feature map; preprocessing, or why it is not applicable] \\
Randomness & [Seeds, repetitions, and fixed versus redrawn reference] \\
Extrapolation & [Fit, range, residual checks and sensitivity for FID$_\infty$] \\
\bottomrule
\end{tabularx}

\Needspace{11\baselineskip}
\subsection{Bench comparison}
\prompt{Score A, B and C under the proposal beside FID and FID$_\infty$.
Identify the protocol for each row. Empty cells below are placeholders, not
zero scores. Explain the direction and units of your proposed score; its scale
need not be comparable to FID's.}
\begin{center}
\begin{tabularx}{\linewidth}{@{}lXXXX@{}}
\toprule
Candidate & Protocol & FID & FID$_\infty$ & Proposed $Q$ \\
\midrule
A & \blank & \blank & \blank & \blank \\
B & \blank & \blank & \blank & \blank \\
C & \blank & \blank & \blank & \blank \\
\bottomrule
\end{tabularx}
\end{center}

\subsection{Estimator bias / Week 5}
\prompt{How does the estimate change with sample size? Separate sampling
variation from bias. State what the measurements support and what remains
unknown; justify any missing evaluation.}

\subsection{Blind spots / Week 6}
\prompt{Describe a construction that challenges your proposal. Distinguish a
failure of the estimator from a limitation of the quantity it estimates.
Explain what happens to C, without assuming that a new metric must detect it.}

\subsection{Correlation with human judgement / Week 10}
\prompt{Evaluate the evidence for agreement with people. A synthetic bench is
not a human study. If you did not run a study, justify that absence and specify
what evidence would be needed instead; do not invent human ratings.}

\subsection{Cost of adoption / Week 11}
\prompt{What would a paper pay in computation, data, human effort and lost
comparability to adopt the proposal? State the assumptions behind the account.}

% Example figure container only; replace with an actual figure and caption.
% For a PDF/PNG figure use: \includegraphics[width=.85\linewidth]{your-figure.pdf}
\begin{figure}[htbp]
  \centering
  \fbox{\parbox[c][18mm][c]{0.88\linewidth}{\centering\small
    [Replace with your diagnostic figure]\par
    No measurements are supplied by this placeholder.}}
  \caption{[Explain the claim this figure tests, the protocol and uncertainty.
  State what the reader can and cannot conclude.]}
  \label{fig:evidence}
\end{figure}

\section{Part C. Where it fails [30 points]}
\prompt{Name the condition under which you would still report FID. State how
badly the proposal fails and which of the four evaluation criteria exposes
that failure. If your answer is ``never'', revisit Part B.}

\subsection{The boundary of the claim}
\prompt{State the strongest claim your evidence supports, then name a claim it
does not support. Explain the practical consequence of that distinction.}

\subsection{Reproducing this report}
\prompt{Point to the separate archive: entry command, seed, environment,
expected outputs and the files behind your tables and figures. Follow the
linked policies for reporting and any tool-use disclosure; this template does
not create a second set of course rules.}

% References are in addition to the six-page report limit. This explicit break
% makes the boundary visible; do not put extra report content in this section.
\clearpage
\section*{References}
\prompt{Replace this prompt with the sources actually used. Cite them where
their ideas enter the report. No references or results have been fabricated
for this starter. You may use BibTeX, biblatex, or the manual example below.}
% \begin{thebibliography}{9}
% \bibitem{your-key} Author. Title. Venue, year.
% \end{thebibliography}
\end{document}
